\section*{Bloque C. Casos con implicancias tarifarias y regulatorias del sector saneamiento}

\noindent\textit{Preguntas 66 a 90. Aplican la teoría del consumidor a la interpretación de estructuras tarifarias, subsidios, costos de servicio, medidas de bienestar y dilemas regulatorios propios de las empresas prestadoras de servicios de saneamiento (EPS).}

\q{Una EPS aplica un cargo fijo mensual de S/ 5 y una tarifa por bloques: de 0 a 10 m$^3$, S/ 1.50 por m$^3$; de 11 a 20 m$^3$, S/ 2.50 por m$^3$; y más de 20 m$^3$, S/ 4.00 por m$^3$. Un usuario consume 25 m$^3$. Su factura mensual es:}
\begin{opts}
    \item S/ 57.50
    \item S/ 65.00
    \item S/ 62.50
    \item S/ 70.00
\end{opts}

\q{En el caso anterior, la tarifa media efectiva por m$^3$ es:}
\begin{opts}
    \item S/ 2.00
    \item S/ 2.50
    \item S/ 4.00
    \item S/ 2.60
\end{opts}

\q{Si el usuario del caso anterior incrementa su consumo en 1 m$^3$, su factura aumenta en:}
\begin{opts}
    \item S/ 4.00
    \item S/ 1.50
    \item S/ 2.50
    \item S/ 5.00
\end{opts}

\q{El cargo fijo de una tarifa de agua se incrementa en S/ 3. Para un usuario con preferencias cuasilineales en agua, $U=v(q)+y$, con consumo interior y ingreso suficiente, el efecto inmediato sobre su consumo de agua es:}
\begin{opts}
    \item aumenta.
    \item disminuye en 3 m$^3$.
    \item nulo.
    \item indeterminado.
\end{opts}

\q{La tarifa por bloques genera una restricción presupuestaria quebrada. Un consumidor puede ubicarse exactamente en el punto de quiebre (el límite del bloque) cuando:}
\begin{opts}
    \item la RMS iguala exactamente al precio marginal del primer bloque.
    \item la RMS queda entre las pendientes de dos bloques adyacentes.
    \item la relación marginal de sustitución es idénticamente nula.
    \item el ingreso monetario disponible del usuario es igual a cero.
\end{opts}

\q{La demanda mensual de agua de un usuario es $Q=20-0.5P$, con $Q$ en m$^3$ y $P$ en S/ por m$^3$. Si $P=10$, la elasticidad precio es:}
\begin{opts}
    \item $-0.25$
    \item $-0.50$
    \item $-2.00$
    \item $-0.33$
\end{opts}

\q{Con la misma demanda, si el precio sube de S/ 10 a S/ 12, el ingreso tarifario del prestador:}
\begin{opts}
    \item aumenta en S/ 18.
    \item cae en S/ 18.
    \item no cambia.
    \item aumenta en S/ 30.
\end{opts}

\q{La evidencia para el consumo urbano de agua suele mostrar demanda inelástica al precio. Si el regulador aprueba un aumento tarifario para financiar inversiones, cabe esperar que:}
\begin{opts}
    \item el consumo caiga en exactamente la misma proporción que el precio.
    \item el ingreso tarifario total disminuya fuertemente tras el alza.
    \item el consumo caiga poco y el ingreso tarifario total aumente.
    \item el consumo aumente de forma sostenida frente al alza tarifaria.
\end{opts}

\q{En la fijación de tarifas Ramsey-Boiteaux para una EPS con costos fijos, el margen sobre el costo marginal debe ser mayor para el grupo de usuarios con demanda:}
\begin{opts}
    \item más elástica.
    \item menos elástica.
    \item idéntica.
    \item perfectamente elástica.
\end{opts}

\q{Una tarifa en dos partes con un cargo fijo igual al excedente del consumidor del usuario más pequeño:}
\begin{opts}
    \item maximiza de forma automática la inclusión de nuevos usuarios.
    \item elimina por completo cualquier pérdida de eficiencia social.
    \item equivale exactamente a fijar un precio igual al costo marginal.
    \item puede excluir a usuarios cuya disposición a pagar no cubre el cargo fijo.
\end{opts}

\q{En un monopolio natural con costo medio decreciente, la tarifa igual al costo medio:}
\begin{opts}
    \item anula el beneficio económico, pero produce menos que el nivel eficiente.
    \item alcanza automáticamente el óptimo social de primer mejor.
    \item genera pérdidas operativas que exigen subsidios fiscales directos.
    \item coincide exactamente con la tarifa lineal igual al costo marginal.
\end{opts}

\q{En el mismo monopolio natural, fijar el precio igual al costo marginal:}
\begin{opts}
    \item genera beneficios extraordinarios positivos para la empresa.
    \item reduce el bienestar social respecto al precio de monopolio.
    \item requiere un subsidio o cargo fijo para cubrir las pérdidas.
    \item elimina por completo la necesidad de regulación económica.
\end{opts}

\q{El subsidio cruzado en saneamiento, mediante el cual los usuarios de mayor consumo financian un descuento a los de menor consumo, se interpreta como:}
\begin{opts}
    \item un impuesto pigouviano destinado a corregir externalidades.
    \item un subsidio cruzado con fin distributivo que distorsiona precios.
    \item una negociación coaseana privada entre usuarios del servicio.
    \item una regla dinámica de precio tope ajustado por productividad.
\end{opts}

\q{Si el consumo básico o de subsistencia es de 10 m$^3$ y se cobra por debajo del costo, el subsidio implícito por usuario (que consume exactamente el bloque) es:}
\begin{opts}
    \item $(precio-costo)\times 15$.
    \item $costo\times 0$ (nulo).
    \item $costo\times 15$.
    \item $(costo-precio)\times 10$.
\end{opts}

\q{Un usuario consume 8 m$^3$ y recibe el bloque de subsistencia. El costo de servicio es S/ 3 por m$^3$ y la tarifa es S/ 1.50 por m$^3$. El subsidio implícito mensual es:}
\begin{opts}
    \item S/ 12
    \item S/ 6
    \item S/ 24
    \item S/ 36
\end{opts}

\q{Si el agua no facturada (ANF) representa 30\% del agua producida y el costo de producción es S/ 2 por m$^3$, el costo efectivo por m$^3$ facturado es aproximadamente:}
\begin{opts}
    \item S/ 1.40
    \item S/ 2.00
    \item S/ 2.86
    \item S/ 3.33
\end{opts}

\q{Si la ANF se reduce de 30\% a 20\%, manteniendo el costo de producción en S/ 2 por m$^3$, el costo efectivo por m$^3$ facturado se reduce en aproximadamente:}
\begin{opts}
    \item S/ 0.14
    \item S/ 0.36
    \item S/ 0.29
    \item S/ 0.50
\end{opts}

\q{Una familia con ingreso mensual de S/ 1,500 recibe una factura de S/ 60. El ratio factura/ingreso es:}
\begin{opts}
    \item 2\%
    \item 5\%
    \item 6\%
    \item 4\%
\end{opts}

\q{Si la disposición a pagar por los primeros 10 m$^3$ es S/ 5 por m$^3$ y la tarifa es S/ 2 por m$^3$, el excedente del consumidor por esos 10 m$^3$ es:}
\begin{opts}
    \item S/ 30
    \item S/ 20
    \item S/ 50
    \item S/ 70
\end{opts}

\q{La variación compensatoria (CV) de un aumento tarifario mide:}
\begin{opts}
    \item el ingreso que habría que quitarle al consumidor, a los precios originales, para alcanzar la nueva utilidad.
    \item solo el cambio en el excedente del consumidor.
    \item el ingreso que habría que dar al consumidor, a los nuevos precios, para restaurar su utilidad original.
    \item la diferencia entre la tarifa media y la tarifa marginal.
\end{opts}

\q{Si la demanda de agua es perfectamente inelástica en el corto plazo, un aumento tarifario:}
\begin{opts}
    \item reduce el consumo en la misma proporción que el incremento.
    \item no reduce el consumo, pero eleva el gasto total de usuarios.
    \item disminuye sustancialmente los ingresos tarifarios de la EPS.
    \item carece de cualquier impacto distributivo sobre los hogares.
\end{opts}

\q{La regulación por precio tope (RPI-X) incentiva a la empresa a:}
\begin{opts}
    \item inflar sus costos operativos para obtener tarifas más altas.
    \item recortar la calidad técnica sin penalización regulatoria.
    \item fijar libremente el precio monopolístico sin restricción.
    \item reducir costos al retener ganancias dentro del período regulado.
\end{opts}

\q{Si el regulador desea maximizar el bienestar social sujeto a que la EPS cubra sus costos sin subsidios, debe aplicar:}
\begin{opts}
    \item precios Ramsey-Boiteaux, con márgenes inversos a la elasticidad.
    \item precios uniformes fijados al nivel estricto del costo marginal.
    \item el precio de monopolio puro sin restricciones de autofinanciamiento.
    \item una tarifa única igual al costo medio sin distinguir demandas.
\end{opts}

\q{En una estimación log-log $\ln Q=a+b\ln P+c\ln Y$, con $b=-0.4$ y $c=0.2$, si el precio sube 10\% y el ingreso 5\%, la cantidad demandada cambia aproximadamente en:}
\begin{opts}
    \item $-1\%$
    \item $+1\%$
    \item $-3\%$
    \item $+3\%$
\end{opts}

\q{Un subsidio focalizado mediante el bloque de consumo básico puede presentar:}
\begin{opts}
    \item ausencia total de errores de inclusión o de exclusión.
    \item filtración a no pobres y subcobertura en pobres de alto consumo.
    \item una progresividad distributiva perfecta en todo tramo.
    \item un costo fiscal neto nulo independientemente del diseño.
\end{opts}
